% Plantilla mínima para TFG (TFG / Trabajo Fin de Grado)
% Archivo: Info_TFG.tex
% Cómo usar: editar título/autor y dividir en capítulos con \include o \input si se desea.
% Compilar (recomendado): usar latexmk o pdflatex + biber/bibtex según bibliografía.
% Ejemplo PowerShell (TeX Live / MiKTeX instalado):
%   latexmk -pdf Info_TFG.tex
% o, paso a paso:
%   pdflatex Info_TFG.tex; biber Info_TFG; pdflatex Info_TFG.tex; pdflatex Info_TFG.tex

\documentclass[12pt,a4paper,oneside]{report}

% Codificación y idioma
\usepackage[T1]{fontenc}
\usepackage[utf8]{inputenc}
\usepackage[spanish]{babel}

% Márgenes y microtipografía
\usepackage[a4paper,margin=2.5cm]{geometry}
\usepackage{microtype}

% Gráficos, tablas y matemáticas
\usepackage{graphicx}
\usepackage{caption}
\usepackage{subcaption}
\usepackage{booktabs}
\usepackage{amsmath,amssymb}

% Hipervínculos
\usepackage[hidelinks]{hyperref}

% Configuración del título/portada (personaliza según la normativa de tu universidad)
\title{Comparación de Nariz Electrónica con olfato de Agentes Caninos para la detección de Drogas}
\author{Raúl Martín-Romo Sánchez\\
	\textit{Grado en Ingeniería Informática del Software}\\
	\textbf{Tutor:} Pablo García Rodríguez\\
    \textbf{Co-tutor:} Jesús Lozano Rogado}
\date{\today}

\begin{document}

% Portada simple
\maketitle

% Capítulos de ejemplo
\chapter{Introducción}
\label{ch:introduccion}
Hablar sobre la idea del proyecto.

\chapter{Motivación}
\label{ch:motivacion}
Hablar sobre la motivación del proyecto (Ayuda al cuerpo canino de la Guardia Civil de Cáceres y admiración por el trabajo de los cuerpos caninos desde pequeño).

\chapter{Planificación}
\label{ch:planificacion}
Planificación del proyecto con metodología Scrum, explicar los diferentes sprints y tareas a realizar añadiendo el diagrama de Gantt.

\chapter{Estado del arte}
\label{ch:estado_arte}
Exponer la inforamación recogida acerca de las narices electrónicas, los sensores y algunos estudios de los diferentes artículos. Además hablar del cuerpo canino de la guardia civil, las funciones que desempeñan y como trabajan.

\chapter{Objetivos}
\label{ch:objetivos}
Describir los objetivos principales.

\chapter{Materiales}
\label{ch:materiales}
Explicar que materiales se han usado y para qué.

\chapter{Desarrollo y entrenamiento de la nariz electrónica}
\label{ch:desarrollo_entrenamiento}
Explicar como ha sido el desarrollo de la e-nose y  como se ha hecho la recogida de datos y el entrenamiento de la nariz electrónica.

\chapter{Pruebas de Comparación, resultados esperados y resultados obtenidos}
\label{ch:pruebas_comparacion}
Describir que tipos de pruebas se han hecho, los resultados esperados, que resultados se han obtenido y si coinciden con los resultados esperados.

\chapter{Discusión}
\label{ch:discusion}
Discusión de los resultados obtenidos.

\chapter{Posibles aplicaciones de la nariz electrónica}
\label{ch:aplicaciones}
Describir las posibles aplicaciones de la nariz electrónica desarrollada como evaluación constante en un sistema de paqueterría o sustitución del cuerpo canino cuando no esten disponibles (están en otro sitio, cansados o enfermos).

\chapter{Futuras mejoras}
\label{ch:futuras_mejoras}
Describir las posibles mejoras como la creación de otra e-nose para detección de explosivos y adapatar ambas narices para que puedan detectar más de una droga/explosivo.

\chapter{Conclusiones}
\label{ch:conclusiones}
Describir las conclusiones del proyecto.

\chapter{Agradecimientos}
\label{ch:agradecimientos}
Agradecer a las personas que han ayudado en el proyecto.

\chapter{Anexos}
\label{ch:anexos}
Material adicional relevante.

\chapter{Referencias}
\label{ch:referencias}
Referencias usadas en el proyecto.

\end{document}

